\section{基础概念总览：语言模型到底在优化什么}

前面的内容已经回答了“为什么要学”和“课程怎么组织”。为了让后面的内容不只是概念堆叠，这里先把语言模型最核心的目标函数讲清楚。
课程里虽然还没有正式进入训练细节，但如果不先知道语言模型到底在优化什么，后面的 tokenization、系统优化、规模律和对齐都很难真正连起来。

一个标准语言模型处理的是长度为 $T$ 的 token 序列 $(x_1,\dots,x_T)$。它的目标不是一次性猜完整段文本，而是学习条件分布
\[
p(x_t \mid x_{<t})
\]
并在整个序列上最小化负对数似然：
\[
\mathcal{L} = - \sum_{t=1}^{T} \log p(x_t \mid x_{<t}).
\]
这就是 next-token prediction 的本质。

\begin{figure}[H]
\centering
\begin{tikzpicture}[node distance=0.95cm, every node/.style={font=\small, align=center}]
\node[draw, rounded corners, fill=blue!10, minimum width=2.8cm, minimum height=0.8cm] (text) {text};
\node[draw, rounded corners, fill=blue!18, right=1.1cm of text, minimum width=2.8cm, minimum height=0.8cm] (tok) {tokenizer};
\node[draw, rounded corners, fill=blue!26, right=1.1cm of tok, minimum width=2.8cm, minimum height=0.8cm] (model) {language model};
\node[draw, rounded corners, fill=blue!34, right=1.1cm of model, minimum width=2.8cm, minimum height=0.8cm] (loss) {cross-entropy};
\draw[-{Latex[length=3mm]}, thick] (text) -- node[above]{encode} (tok);
\draw[-{Latex[length=3mm]}, thick] (tok) -- (model);
\draw[-{Latex[length=3mm]}, thick] (model) -- node[above]{predict next token} (loss);
\end{tikzpicture}
\caption{语言模型最小闭环：原始文本经过 tokenization，再进入模型，最后用 next-token loss 训练}
\end{figure}

\begin{table}[H]
\centering
\begin{tabular}{p{3cm}p{4.2cm}p{5.8cm}}
\toprule
\textbf{符号} & \textbf{含义} & \textbf{为什么重要} \\
\midrule
$x_t$ & 第 $t$ 个 token & 模型的基本预测单位 \\
$x_{<t}$ & 前文上下文 & 决定条件概率 \\
$p(x_t|x_{<t})$ & 预测分布 & 训练和评估都围绕它展开 \\
$\mathcal{L}$ & 负对数似然 / cross-entropy & 真正被优化的目标 \\
\bottomrule
\end{tabular}
\caption{本门课后面会反复使用的基础符号}
\end{table}

\begin{knowledgebox}{为什么 next-token objective 很关键}
它有三个好处：第一，训练信号极其密集，几乎每个 token 都能提供监督；第二，它天然适配大规模文本；第三，很多看似复杂的能力都能在这个目标下涌现出来，所以它既简单又足够强。
\end{knowledgebox}

\begin{importantbox}{一个容易忽略的点}
“语言模型”不是在生成“句子”，而是在估计 token 序列分布。句子只是人类解释模型输出的方式，训练目标本身更原子、更机械。
\end{importantbox}

\subsection{从 loss 到 perplexity}

课程里经常会遇到 perplexity 这个词。它不是神秘指标，本质上就是把平均负对数似然重新指数化，便于理解“模型对下一个 token 的平均惊讶程度”。
如果把交叉熵记作 $H$，那么 perplexity 近似就是
\[
\mathrm{PPL} = \exp(H).
\]
因此，降低 perplexity 和降低交叉熵是同一件事，只是表达方式不同。

\begin{warningbox}{不要把 perplexity 当成万能指标}
perplexity 能反映语言建模能力，但不直接等价于“更有用”或“更安全”。后面的数据和对齐章节会不断提醒你这一点。
\end{warningbox}

